%% kapa: exam=1 round=27 year=2016 period=2 form=A subject=appraisal_law subject_ko=감정평가관계법규 seq=8 number=8
%% source: Q-Net 감정평가사 기출문제 (cst00309) artlSeq=5169402 file=extracted/1008/5169402/제27회 감정평가사 1차 문제/제27회 감정평가사 2교시문제(A형).pdf
%% conversion: kordoc 4.24.51 -> tools/gen1.py (auto); review=proofread-sonnet; flags=none
\begin{question}{8}
국토의 계획 및 이용에 관한 법령상 개발밀도관리구역에 관한 설명으로 옳지 않은 것은?
\begin{choices}
  \item 주거ㆍ상업 또는 공업지역에서의 개발행위로 기반시설의 처리능력이 부족할 것이 예상되는 지역 중 기반시설의 설치가 곤란한 지역을 개발밀도관리구역으로 지정할 수 있다.
  \item 개발밀도관리구역을 지정할 때 개발밀도관리구역에서는 당해 용도지역에 적용되는 건폐율 또는 용적률을 강화하여 적용한다.
  \item 개발밀도관리구역을 지정하기 위해서는 지방도시계획위원회의 심의를 거쳐야 한다.
  \item 개발밀도관리구역의 경계는 특색 있는 지형지물을 이용하는 등 경계선이 분명하게 구분되도록 하여야 한다.
  \item 개발밀도관리구역의 지정기준을 정할 때 고려되는 기반시설에 학교는 포함되지 않는다.
\end{choices}
\end{question}
\answer{5}
